\documentclass[../../main.tex]{subfiles}
\graphicspath{{../../image/}{../../figure/}}
\begin{document}

\section{连续 Glimm 泛函、稀疏波的展开与解的结构}\label{section:ch14-连续glimm泛函与稀疏波展开}

本节借助前沿追踪算法，说明\cref{theorem:ch14-唯一性}建立的解的唯一性在刻画标准 Riemann 半群轨道的结构方面具有有用的含义。特别地，将推导出解中弱激波与稀疏前沿的渐近展开。

出发点是构造一列 $U$，其初值 $U_0$ 位于 $\mathscr{D}$ 中，且其像 $U(\cdot,t)=S_tU_0(\cdot)$ 由\eqref{align:ch13-sup-bound}给出。然而，通过把激波与稀疏波视为连续曲线而非前沿追踪近似中的折线，我们采用稍微不同的观点。因而，与\cref{section:ch14-真非线性系统的前沿追踪}中近似解的情形不同，允许稀疏波以连续扇形出现。



为下面使用，回顾\eqref{align:ch09-9-3-9}中定义的 Riemann 不变量 $w_i$，它们在 $i$-特征曲线上为常数。



对解 $U$ 的任一跳跃间断 $x=y(t)$，具有单侧极限 $U_\pm=U(y(t)\pm,t)$，通过在 $U_-$ 与 $U_+$ 之间求解经典 Riemann 问题分解跳跃间断，并测量出射波扇中所有 $i$-波的总强度 $|\varepsilon_i|$，来定义其\textbf{强度} $|\gamma_i|$（$i=1,\cdots,n$）。当解由单 $i$-波组成时，强度 $|\gamma_i|$ 简化为激波振幅的绝对值或相应稀疏波的自然强度。

我们现通过连续泛函建立 $U$ 的稳定性，这些泛函监控波型的振荡与相互作用，在形式上类似于\eqref{align:ch13-Q-change}中定义的不连续泛函 $L$ 与 $Q$。



对 $i=1,\cdots,n$，令 $\mathscr{S}_i(t)$ 表示穿过 $t$-时间线的所有 $i$-波集合，并令 $\mathscr{C}_i(t)$ 表示相应 $i$-特征曲线（包括由解生成的 $i$-特征）的集合。定义
\begin{align}
\Theta_i(t)=\sum_{\gamma\in\mathscr{S}_i(t)}|\gamma_i|,\qquad \Xi_i(t)=\sup_{C\in\mathscr{C}_i(t)}\operatorname{osc}_C w_i.
\label{align:ch14-ThetaXi}
\end{align}
$\Theta_i(t)$ 穿过 $t$-时间线测量 $i$-波的总强度，而 $\Xi_i(t)$ 沿任一 $i$-特征测量 Riemann 不变量 $w_i$ 的最大振荡。

两个不同波 $\alpha\in\mathscr{S}_i(t)$ 与 $\beta\in\mathscr{S}_j(t)$，若它们穿过 $t$-时间线时满足 $x_\alpha<x_\beta$，且 $i>j$ 或 $i=j$ 且至少一个波是激波，则称它们是\textbf{接近的}。对 $i=1,\cdots,n$，定义
\begin{align}
Q_i(t)=\sum_{\substack{\alpha\in\mathscr{S}_i(t)\\ \beta\in\bigcup_{k\neq i}\mathscr{S}_k(t)\\ \alpha,\beta\text{ 接近}}}|\alpha_i||\beta_k|.
\label{align:ch14-Qi}
\end{align}

\begin{lemma}\label{lemma:ch14-连续glimm泛函估计}
对任一 $t\in[0,\infty)$，
\begin{align}
\Theta_i(t)-\Xi_i(t)\leqslant\Theta_i(0)-\Xi_i(0)+c_{17}\{Q_i(0)-Q_i(t)\},
\label{align:ch14-ThetaXi-est}
\end{align}
其中 $c_{17}$ 是仅依赖于 $F$ 及其导数界的常数。
\end{lemma}

\begin{proof}
为证明\eqref{align:ch14-ThetaXi-est}，须考察当时间增长时 $\Theta_i(t)-\Xi_i(t)$ 的变化。首先，当波碰撞发生时，穿过碰撞点的出射波的总强度由\cref{theorem:ch09-波扇相互作用}控制。当碰撞涉及两个不同族 $i$-波与 $j$-波（$i<j$）时，出射 $i$-波的总强度仅以 $O(1)|\alpha_i||\beta_j|$ 改变。另一方面，$Q_i(t)$ 以约 $|\alpha_i||\beta_j|$ 的量下降，这是因为这些波在碰撞前是接近的。类似论证适用于两个 $i$-波碰撞的情形。



当 $i$-波与 $j$-波（$j\neq i$）碰撞时，$\Xi_i(t)$ 也可能改变。事实上，沿 $i$-特征的 Riemann 不变量 $w_i$ 的振荡受与 $j$-波的相互作用影响。$Q_i(t)$ 中记录的 $i$-波与 $j$-波的接近波对提供了足够的抵消。因而在每种情形下，$\Theta_i(t)-\Xi_i(t)$ 的增量都由 $Q_i(t)$ 的下降所主导，这给出\eqref{align:ch14-ThetaXi-est}。
\end{proof}

特别地，\cref{lemma:ch14-连续glimm泛函估计}蕴含 $Q_i(t)$ 是非增的，因而波相互作用总量有界。

作为\cref{lemma:ch14-连续glimm泛函估计}的推论，我们现在刻画解的结构。

\begin{theorem}\label{theorem:ch14-稀疏波展开}
设 $U$ 是 $S_t$ 的轨道，从 $U_0\in\mathscr{D}$ 发出。则对任一 $i=1,\cdots,n$，$i$-稀疏波的强度随时间展开：
\begin{align}
\sum_{\text{所有 }i\text{-稀疏波}}|\gamma_i|\leqslant\frac{c_{18}}{t}.
\label{align:ch14-rarefaction-spread}
\end{align}
特别地，当 $t\to\infty$ 时，任何 $i$-稀疏波的强度趋于零。
\end{theorem}

\begin{proof}
考虑任一 $i$-稀疏波扇，其左端状态 $U_l$ 与右端状态 $U_r$。由于 $w_i$ 在稀疏波扇上严格单调，$w_i(U_r)-w_i(U_l)$ 以稀疏波强度为界。$w_i$ 的变化沿 $i$-特征累积，而沿每条 $i$-特征，$w_i$ 的振荡被\eqref{align:ch14-ThetaXi-est}控制。

通过波相互作用分析可得 $i$-稀疏波强度的 $1/t$ 衰减。其机制是：弱激波与稀疏波的相互作用趋于把稀疏波``拉开''，使其强度随时间扩散。更精确地，稀疏波在时刻 $t$ 的宽度以其初强度与相互作用时间的乘积度量，而其强度（高度）以宽度为界。因而当 $t$ 增长时，稀疏波展开，强度按 $1/t$ 衰减。这给出\eqref{align:ch14-rarefaction-spread}。
\end{proof}

\begin{theorem}\label{theorem:ch14-解的渐近结构}
当 $t\to\infty$ 时，解 $U(\cdot,t)=S_tU_0(\cdot)$ 渐近趋于由激波接触间断分离的常状态配置。特别地，在任一固定有界区域中，$U(\cdot,t)$ 收敛到某个阶梯函数，其跳跃位置以激波速度传播。
\end{theorem}

\begin{proof}
由\cref{theorem:ch14-稀疏波展开}，稀疏波强度随 $1/t$ 衰减。同时，$Q_i(t)$ 的非增性蕴含波相互作用总量有界，因而碰撞次数随时间减少。当所有稀疏波消失后，解简化为由激波与接触间断组成的分段常值函数，这些波以特征速度或激波速度直线传播。因而在任何固定有界区域中，$U(\cdot,t)$ 最终收敛到某阶梯函数。
\end{proof}

\begin{theorem}\label{theorem:ch14-弱激波稳定性}
在 $U$ 的演化中产生的弱激波的强度随时间不增，且任意弱激波在渐近极限中可任意小。
\end{theorem}

\begin{proof}
弱激波仅在波碰撞中生成，而波相互作用总量由 $Q_i(t)$ 的非增性有界。此外，跨任何碰撞，出射弱激波的强度以入射波强度的乘积为界。因而弱激波强度随时间不增，且最终趋于零。
\end{proof}

\end{document}
